\section*{TÓM TẮT ĐỒ ÁN}
\phantomsection\addcontentsline{toc}{section}{\numberline {}TÓM TẮT ĐỒ ÁN}

\textbf{Tiếng Việt:}

Đồ án này đánh giá hiệu quả cơ chế Hybrid Automatic Repeat reQuest (HARQ) trong mạng 5G Phi Mặt Đất (Non-Terrestrial Network -- NTN) bằng phương pháp mô phỏng Monte Carlo. Mô hình kênh Rician phẳng đặc trưng cho liên kết vệ tinh -- mặt đất được xây dựng theo TR 38.811, với hệ số K = 15 dB và bốn quỹ đạo tiêu biểu: LEO 600 km, LEO 1200 km, MEO 10000 km và GEO 35786 km. Hai sơ đồ kết hợp HARQ được so sánh -- Chase Combining (CC) và Incremental Redundancy (IR) -- với mô hình ngưỡng thông tin tương hỗ (MI-threshold) sử dụng khoảng cách thực thi LDPC 1{,}5 dB. Kết quả cho thấy IR-HARQ với bốn lần phát mang lại lợi ích 8{,}3 dB về SNR so với không có HARQ tại BLER = $10^{-3}$ (MCS A, QPSK $r = 1/2$). Tuy nhiên, ràng buộc số tiến trình HARQ tối đa là 32 (TS 38.214) gây ra hiện tượng pipeline stall nghiêm trọng tại MEO và GEO, nơi N$_{\min}$ vượt xa 32. Tại GEO, pipeline stall giảm hiệu suất phổ tần xuống còn 5{,}89\% so với lý thuyết, khiến giải pháp vô hiệu hóa HARQ kết hợp RLC ARQ tăng SE Goodput lên 16{,}7 lần. Đồ án cũng trình bày phân tích năng lượng/bit theo quỹ đạo và khuyến nghị cấu hình HARQ tối ưu cho từng quỹ đạo.

\vspace{12pt}
\textbf{Tiếng Anh (Abstract):}

This thesis evaluates the effectiveness of the Hybrid Automatic Repeat reQuest (HARQ) mechanism in 5G Non-Terrestrial Networks (NTN) through Monte Carlo simulation. A Rician flat-fading channel model characterising the satellite-to-ground link is built in accordance with TR 38.811, with K-factor of 15 dB over four representative orbits: LEO 600 km, LEO 1200 km, MEO 10000 km, and GEO 35786 km. Two HARQ combining schemes -- Chase Combining (CC) and Incremental Redundancy (IR) -- are compared using a mutual-information threshold model with a 1.5 dB LDPC implementation gap. Results show that IR-HARQ with four transmissions yields an 8.3 dB SNR gain over no-HARQ at BLER $= 10^{-3}$ (MCS A, QPSK $r = 1/2$). However, the 3GPP constraint of at most 32 HARQ processes (TS 38.214) causes severe pipeline stall at MEO and GEO orbits where $N_{\min}$ far exceeds 32. At GEO, this stall reduces spectral efficiency to 5.89\% of the theoretical maximum, making HARQ-disabled (RLC ARQ) operation 16.7$\times$ more efficient in SE Goodput. An energy-per-bit analysis across orbits and recommended HARQ configurations per orbit are also presented.

\cleardoublepage
